\documentclass[11pt,a4paper,openany]{book}
\input{userguide-preamble}

\hypersetup{
  pdftitle={Snapdragon Games — Intro for coders},
  pdfauthor={Snapdragon Systems},
  pdfsubject={Entry-level orientation for the Snapdragon Games portal},
}

\begin{document}

% ── Cover ─────────────────────────────────────────────────────────────────
\begin{titlepage}
\thispagestyle{empty}
\noindent\begin{tikzpicture}[remember picture,overlay]
  \fill[snapBg] (current page.south west) rectangle (current page.north east);
  \fill[snapAccent,opacity=0.18] ([yshift=-0\paperheight]current page.north west) circle (8cm);
  \fill[snapAccentTwo,opacity=0.10] ([xshift=4cm,yshift=-2cm]current page.north west) circle (5cm);
\end{tikzpicture}
\vspace*{4cm}\noindent
{\sffamily\bfseries
  {\color{snapAccentTwo}\fontsize{18}{22}\selectfont SNAPDRAGON GAMES}\\[3mm]
  {\color{white}\fontsize{42}{50}\selectfont Intro for coders}\par}
\vspace{1.5cm}
\noindent\color{white!75!snapAccentTwo}\rule{0.4\linewidth}{1.2pt}
\vspace{0.8cm}
\noindent
{\sffamily\color{white!85}\large
A 12-section orientation to the federated portal: where code lives,
how a click reaches a frame, and the small set of patterns that hold
the stack together.}
\vfill\noindent
{\sffamily\color{white!70}\small
User guide companion to \texttt{intro-for-coders.pptx}.\\
Generated: \snapbuilddate{} \snapbuildtime}
\end{titlepage}

\frontmatter
\pagestyle{empty}
\tableofcontents
\clearpage
\mainmatter
\pagestyle{fancy}

% ── Chapter 1 ─────────────────────────────────────────────────────────────
\chapter{What is Snapdragon Games?}
A browser-native portal that hosts multiple Godot 4 games on one page.
The Angular shell wraps each game with HUD, login, admin; a Spring Boot
BFF streams game state and controls over Server-Sent Events. PostgreSQL
holds the catalog and player state, Cassandra is the cross-pod movement
cache and audit log, and Casdoor handles OAuth2/OIDC behind a federated
login remote.

The portal is delivered as one \texttt{docker compose} stack of
roughly twelve containers: a host shell, four federated route or game
remotes, an asset-tier nginx for the WASM bundles, the BFF, two
databases, the identity provider, and a few init containers.

\section{Three sentences}
The host shell is a native-federation Angular host that loads its
\texttt{login}, \texttt{admin}, and per-game routes from sibling bundles
served by per-remote nginx containers. The BFF runs on Spring Boot 3.4
WebFlux + R2DBC + JDBC, exposes the SSE stream, the federation manifest,
the upload pipeline, and the per-game launch / access endpoints, and
persists state in Postgres with a Cassandra cache. Each Godot game
lives in its own project under \snapfile{games/<slug>/}, exports a
self-contained WASM bundle, and is delivered to the browser through
the dedicated \texttt{wasm-fragment} nginx remote with the COOP/COEP/CORP
headers WASM games need.

\section{Three-tier topology}
\snapdiagram{topology.tex}

% ── Chapter 2 ─────────────────────────────────────────────────────────────
\chapter{Where the code lives}
\begin{itemize}
  \item \snapfile{frontend/} — Angular host (one project).
  \item \snapfile{frontend/projects/} — the four federated remotes
        (\texttt{login}, \texttt{admin}, \texttt{games-three-wizards},
        \texttt{games-dodecahedron}).
  \item \snapfile{wasm-fragment/} — asset-tier nginx that serves Godot
        bundles.
  \item \snapfile{backend/} — Spring Boot BFF (WebFlux + R2DBC).
  \item \snapfile{games/<slug>/} — one Godot project per game.
  \item \snapfile{extensions/snapdragon-bridge/} — reusable Rust
        GDExtension (native + WASM).
  \item \snapfile{infra/}, \snapfile{casdoor/}, \snapfile{e2e/} — init
        containers, identity config, end-to-end tests.
\end{itemize}

% ── Chapter 3 ─────────────────────────────────────────────────────────────
\chapter{How a player click reaches a frame}
\snapdiagram{runtime-flow.tex}

\begin{enumerate}
  \item The host's \snapfile{HomeComponent} reads \texttt{SCENE\_TYPE\_TO\_GAME\_SLUG}
        and forks the renderer based on the active level's scene type.
  \item Federated games \texttt{POST /api/games/\{slug\}/launch} to
        mint a 256-bit URL-safe \texttt{launch\_uid}.
  \item \texttt{RemoteLoaderService} loads the matching per-game remote
        through \texttt{globalThis.\_\_SNAPDRAGON\_LOAD\_REMOTE\_MODULE\_\_}.
  \item The remote mounts a Web Component that builds the iframe URL.
  \item \texttt{wasm-fragment} serves \texttt{index.html} with the
        COOP/COEP/CORP headers WASM bundles need.
  \item Godot pulls scene JSON from
        \texttt{GET /api/levels/\{id\}/scene} and renders.
\end{enumerate}

% ── Chapter 4 ─────────────────────────────────────────────────────────────
\chapter{Server-Sent Events: one stream, five channels}
\snapdiagram{game-stream-controller.tex}
\snapmarblelegend{
  \snapmarbleshape{marbleFlux} & \texttt{PlayerMovement}
    & \texttt{characterId} {\rm $\cdot$} \texttt{x,y,z} {\rm $\cdot$} \texttt{rotationY} {\rm $\cdot$} \texttt{action} {\rm $\cdot$} \texttt{pressed} {\rm $\cdot$} \texttt{tick}. \\
  \snapmarbleshape{marbleFlux} & \texttt{WorldEvent}
    & \texttt{type} (\texttt{spawn}/\texttt{despawn}/\texttt{transition}) {\rm $\cdot$} \texttt{payload}. \\
  \snapmarbleshape{marbleFlux} & \texttt{CharacterEvent}
    & \texttt{characterId} {\rm $\cdot$} \texttt{hp}, \texttt{mp}, \texttt{xp} {\rm $\cdot$} flags. \\
  \snapmarbleshape{marbleFlux} & \texttt{GameTimeTick}
    & \texttt{day} {\rm $\cdot$} \texttt{hour} {\rm $\cdot$} \texttt{minute} {\rm $\cdot$} \texttt{phase}. \\
  \snapmarbleshape{marbleSig}  & ping (15-second keepalive)
    & no payload — fires immediately to flush \texttt{onopen}. \\
  \snapmarbleshape{marbleSSE}  & \texttt{ServerSentEvent}
    & merged: \texttt{event}, \texttt{data}, \texttt{id}. \\
}

\begin{snapcaution}[title={The autoCancel trap}]
\texttt{Sinks.many().multicast().onBackpressureBuffer(N)} defaults to
\texttt{autoCancel=true}: when every subscriber disconnects (iframe
remount, EventSource reconnect, network blip), the sink terminates
permanently and all subsequent \texttt{tryEmitNext} calls silently
drop. Every sink in this project passes \texttt{autoCancel=false}
explicitly. If a sink ``works once and then dies'', this is the cause.
\end{snapcaution}

% ── Chapter 5 ─────────────────────────────────────────────────────────────
\chapter{How to run it}
\begin{itemize}
  \item \texttt{make build} — full source build (rust + godot + backend + frontend).
  \item \texttt{docker compose up -d} — bring the whole stack up.
  \item \texttt{make e2e} — full Playwright suite against the live stack.
  \item \texttt{make e2e-native} — same, against the GraalVM native binary.
\end{itemize}
Open \texttt{http://localhost:8080}; the home page shows the visible
games for the current Casdoor user (or all games if you're admin).

\begin{figure}[!h]
\centering
\includegraphics[width=0.95\linewidth]{home-page.png}
\caption{The portal home page — federated host shell with cards for
each game the current user can see.}
\end{figure}

% ── Chapter 6 ─────────────────────────────────────────────────────────────
\chapter{Your first change}
\begin{itemize}
  \item Tweaking a controller / DTO / SQL? $\rightarrow$
        \snapfile{backend/} + a Liquibase changeset.
  \item Tweaking a route / login UI / admin? $\rightarrow$
        \snapfile{frontend/projects/<remote>/}.
  \item Tweaking the game? $\rightarrow$ \snapfile{games/<slug>/} in
        the Godot editor (after \texttt{make build-rust}).
  \item Tweaking the runtime data flow? $\rightarrow$
        \snapfile{HomeComponent} + \snapfile{RemoteLoaderService}.
\end{itemize}
Always run the matching tests: \texttt{make test-backend},
\texttt{make test-frontend}.

% ── Chapter 7 ─────────────────────────────────────────────────────────────
\chapter{Roles you'll meet}
The Casdoor seed data ships twelve users. Every one can log in to the
catalog out of the box (pinned by
\reposrc{e2e/features/auth/per-role-login.feature}); what differs is
what each one sees and is allowed to do after that.

\begin{itemize}
  \item \texttt{player}, \texttt{player1}, \texttt{player2},
        \texttt{player3} --- ordinary players. Land on the catalog;
        can create + join sessions; cannot reach \texttt{/admin}.
  \item \texttt{admin} --- ROLE\_ADMIN. Sees every game in the catalog
        regardless of \texttt{game\_access}; can grant mod / publisher
        scope; observes any live session.
  \item \texttt{mod} --- ROLE\_MOD with no per-game scope. Authenticates
        but every mod-tier action returns 403 until an admin grants a
        \texttt{mod\_game\_scope} row.
  \item \texttt{mod-three-wizards}, \texttt{mod-dodecahedron} --- mod
        with a specific game scope. Can act on that game's sessions;
        gets 403 on any other game.
  \item \texttt{publisher} --- ROLE\_PUBLISHER with no scope. Same
        shape as \texttt{mod}: authenticates, but any
        \texttt{/api/publisher/*} call needs a matching
        \texttt{publisher\_game\_scope} row.
  \item \texttt{publisher-three-wizards}, \texttt{publisher-dodecahedron},
        \texttt{publisher-flatland} --- per-game publisher fixtures.
        The Flatland one auto-seeds an RPG-maker + VN-maker project
        owned by that user, so logging in as them opens the maker UI
        on a ready-to-build project tree.
\end{itemize}

The two-axis split (\emph{role} in Casdoor, \emph{game scope} in
Postgres) is the same shape for mod and publisher --- decision 015 +
018 in the wiki cover the wiring.

% ── Chapter 8 ─────────────────────────────────────────────────────────────
\chapter{Multiplayer in one paragraph}
Two players in the \emph{same} session share state via per-session SSE
sinks. \texttt{GameStateService.applyControl(event, characterId,
sessionId)} fans every event to the matching
\texttt{ConcurrentHashMap<String, SessionSinks>} entry; subscribers on
\texttt{/api/game/stream?sessionId=$\alpha$} only see events for
session $\alpha$. Two players in \emph{different} sessions never see
each other's events --- the isolation contract is pinned by
\snapfile{SessionIsolationTest} on the BFF side and
\reposrc{e2e/features/multiplayer/rpg-two-players.feature} on the
browser side.

Visibility comes from two more layers. The lobby filters sessions by
the caller's \texttt{user\_relationships} graph: a \texttt{friends}
session is invisible to non-friends; a \texttt{blocked} user is
invisible to the blocker in both directions (lobby list AND the
in-game player roster). The friend-graph happy path is pinned by
\reposrc{e2e/features/social/friends-and-blocks.feature}.

% ── Chapter 9 ─────────────────────────────────────────────────────────────
\chapter{Reading list}
\begin{itemize}
  \item \reposrc{README.md} — installation, build, deploy, troubleshooting.
  \item \reposrc{docs/PLAN-module-federation-and-native.md} — the
        eight-section plan that produced the current shape.
  \item \reposrc{docs/GODOT_DEVELOPER_GUIDE.md} — slug contract,
        presets, upload flow.
  \item \reposrc{docs/publishing/architecture/build/architecture.pdf} —
        full typeset architecture reference.
  \item \reposrc{docs/publishing/javadoc-output/index.html} — branded
        Javadoc with marble-style references.
\end{itemize}

\end{document}
